\vspace{0.5em}\noindent\textbf{Tạo S3 bucket}\par

\begin{enumerate}
\def\labelenumi{\arabic{enumi}.}
\item
  Đi đến S3 management console
\item
  Trong Bucket console, chọn \textbf{Create bucket}
\item
  Trong Create bucket console
\end{enumerate}

\begin{itemize}
\item
  Đặt tên bucket: chọn 1 tên mà không bị trùng trong phạm vi toàn cầu
  (gợi ý: lab\textless số-lab\textgreater\textless tên-bạn\textgreater)
\item
  Giữ nguyên giá trị của các fields khác (default)
\item
  Kéo chuột xuống và chọn \textbf{Create bucket}
\item
  Tạo thành công S3 bucket
\end{itemize}

\vspace{0.5em}\noindent\textbf{Kết nối với EC2 bằng session
manager}\par

\begin{itemize}
\item
  Trong workshop này, bạn sẽ dùng AWS Session Manager để kết nối đến các
  EC2 instances. Session Manager là 1 tính năng trong dịch vụ Systems
  Manager được quản lý hoàn toàn bởi AWS. System manager cho phép bạn
  quản lý Amazon EC2 instances và các máy ảo on-premises (VMs)thông qua
  1 browser-based shell. Session Manager cung cấp khả năng quản lý phiên
  bản an toàn và có thể kiểm tra mà không cần mở cổng vào, duy trì máy
  chủ bastion host hoặc quản lý khóa SSH.
\item
  First Cloud AI Journey
  \href{https://000058.awsstudygroup.com/1-introduce/}{Lab} để hiểu sâu
  hơn về Session manager.
\end{itemize}

\begin{enumerate}
\def\labelenumi{\arabic{enumi}.}
\item
  Trong AWS Management Console, gõ Systems Manager trong ô tìm kiếm và
  nhấn Enter:
\item
  Từ \textbf{Systems Manager} menu, tìm \textbf{Node Management} ở thanh
  bên trái và chọn \textbf{Session Manager}:
\item
  Click Start Session, và chọn EC2 instance tên
  \textbf{Test-Gateway-Endpoint}.
\end{enumerate}

Phiên bản EC2 này đã chạy trong ``VPC cloud'' và sẽ được dùng để kiểm
tra khả năng kết nối với Amazon S3 thông qua điểm cuối Cổng mà bạn vừa
tạo (s3-gwe).

Session Manager sẽ mở browser tab mới với shell prompt: sh-4.2 \$

Bạn đã bắt đầu phiên kết nối đến EC2 trong VPC Cloud thành công. Trong
bước tiếp theo, chúng ta sẽ tạo một S3 bucket và một tệp trong đó.
\#\#\#\# Create a file and upload to s3 bucket

\begin{enumerate}
\def\labelenumi{\arabic{enumi}.}
\item
  Đổi về ssm-user's thư mục bằng lệnh ``cd \textasciitilde{}''
\item
  Tạo 1 file để kiểm tra bằng lệnh ``fallocate -l 1G testfile.xyz'', 1
  file tên ``testfile.xyz'' có kích thước 1GB sẽ được tạo.
\item
  Tải file mình vừa tạo lên S3 với lệnh ``aws s3 cp testfile.xyz
  s3://your-bucket-name''. Thay your-bucket-name bằng tên S3 bạn đã tạo.
\end{enumerate}

Bạn đã tải thành công tệp lên bộ chứa S3 của mình. Bây giờ bạn có thể
kết thúc session.

\vspace{0.5em}\noindent\textbf{Kiểm tra object trong S3
bucket}\par

\begin{enumerate}
\def\labelenumi{\arabic{enumi}.}
\tightlist
\item
  Đi đến S3 console.\\
\item
  Click tên s3 bucket của bạn
\item
  Trong Bucket console, bạn sẽ thấy tệp bạn đã tải lên S3 bucket của
  mình
\end{enumerate}

\vspace{0.5em}\noindent\textbf{Tóm tắt}\par

Chúc mừng bạn đã hoàn thành truy cập S3 từ VPC. Trong phần này, bạn đã
tạo gateway endpoint cho Amazon S3 và sử dụng AWS CLI để tải file lên.
Quá trình tải lên hoạt động vì gateway endpoint cho phép giao tiếp với
S3 mà không cần Internet gateway gắn vào ``VPC Cloud''. Điều này thể
hiện chức năng của gateway endpoint như một đường dẫn an toàn đến S3 mà
không cần đi qua pub lic Internet.
